% Generated by report/make_restated.py from theory/A1_lower_bound.tex: do not edit (edit the note and rerun).
\begin{restated}{Definition}{def:ifo}{A1}{0}{incremental first-order oracle}
Following \cite{agarwal2015lower} (Definition~2), the \emph{incremental first-order oracle} (IFO) takes as input a point $\theta\in\R^d$ and an index $i\in\{1,\dots,n\}$ and returns the pair $(f_i(\theta),\fip(\theta))$. An IFO algorithm is an algorithm whose specification depends on $F$ only through such calls; its cost is the number of calls. The \emph{proximal} IFO (PIFO) of \cite{woodworth2016tight,han2024lower} additionally returns $\operatorname{prox}_{f_i}(\theta,\beta)=\argmin_u\{f_i(u)+\frac\beta2\norm{u-\theta}^2\}$: for any $\beta>0$ in \cite{han2024lower}, who write $\frac1{2\gamma}$, $\gamma>0$, for $\frac\beta2$; in display (3) of \cite{woodworth2016tight} the $\argmin$ is over the domain $\mathcal X$ and $\beta=0$ is allowed (``one prox oracle call (with $\beta=0$) is sufficient for exact optimization'' of a single function). The oracle of \cite{woodworth2016tight} is thus the stronger one, and lower bounds proved for it hold a fortiori for the other.
\end{restated}
